% Solution: Tadelis (2013) Exercise 7.10 "Perfect and Imperfect Recall". Input by \solitem.
% New in the 30-item list. Check: check_B_L03.py.
\begin{exsol}{Tadelis Exercise 7.10}{150}{Exam}{4}
\asked In the game of Figure~7.14: (a) pure-strategy sets; (b) every behavioral strategy of
player~1 has an outcome-equivalent mixed strategy; (c) the same for player~2; (d) whether (c)
fails once player~2's two bottom information sets are merged.

\method Pure, mixed and behavioral strategies (Lecture~3, slides~25 and~27; Tadelis \S7.2.1--7.2.2,
Defs.~7.4--7.6, pp.~137--140); perfect recall and Kuhn's equivalence theorem (Lecture~3, slide~28;
Tadelis Def.~7.7 and the discussion on pp.~141--142). Outcome equivalence is checked node by node:
two strategies of a player are equivalent if, against every pure strategy of the opponent (and
then, by linearity, every mixed one), they put the same probability on each terminal node.

\solution
\textbf{The game.} Player~2 moves first ($a$ or $b$). After $a$, player~1 chooses $L$ or $R$;
after $b$, he chooses $W$ or $E$. Player~2 then moves again without seeing player~1's choice: at
information set $h_a = \{aL, aR\}$ with actions $t, z$, and at $h_b = \{bW, bE\}$ with actions
$x, y$. Payoffs $(u_1,u_2)$: $(2,1)$ after $Lt$ and $Wx$, $(1,2)$ after $Rz$ and $Ey$, $(0,0)$
otherwise.

\textbf{(a) Pure strategies.} Player~1 has two information sets (after $a$, after $b$), each with
two actions:
\[
  S_1 = \{L, R\} \times \{W, E\} = \{LW, LE, RW, RE\} \quad (4).
\]
Player~2 has three information sets (the root, $h_a$, $h_b$):
\[
  S_2 = \{a, b\} \times \{t, z\} \times \{x, y\} \quad (2^3 = 8),
\]
e.g.\ $(a, t, x)$, $(a, t, y)$, \dots, $(b, z, y)$.

\textbf{(b) Player 1: behavioral $\Rightarrow$ mixed.} A behavioral strategy is $\beta_1 = (\lambda,
\omega)$, with $\lambda = \Pr(L)$ at the left node and $\omega = \Pr(W)$ at the right node. Define
the mixed strategy as the product
\[
  \sigma_1(LW) = \lambda\omega,\quad \sigma_1(LE) = \lambda(1-\omega),\quad
  \sigma_1(RW) = (1-\lambda)\omega,\quad \sigma_1(RE) = (1-\lambda)(1-\omega).
\]
Fix any pure $s_2 = (r, m_a, m_b)$. If $r = a$, the play reaches player~1's left node and ends at
$(a, L, m_a)$ or $(a, R, m_a)$. Under $\beta_1$ these have probabilities $\lambda$ and $1-\lambda$.
Under $\sigma_1$, $(a, L, m_a)$ occurs iff the drawn plan starts with $L$:
\[
  \sigma_1(LW) + \sigma_1(LE) = \lambda\omega + \lambda(1-\omega) = \lambda,
\]
and $(a, R, m_a)$ has probability $1 - \lambda$. The case $r = b$ is symmetric with $\omega$. All other
terminal nodes have probability 0 under both. So $\beta_1$ and $\sigma_1$ induce the same distribution
against every pure $s_2$, hence (expected values are linear in $\sigma_2$) against every mixed
$\sigma_2$ too. Since player~1 is called once per play, only the marginals matter: any $\sigma_1$ with
$\sigma_1(LW) = \mu$, $\sigma_1(LE) = \lambda - \mu$, $\sigma_1(RW) = \omega - \mu$,
$\sigma_1(RE) = 1 - \lambda - \omega + \mu$, $\mu \in [\max\{0, \lambda+\omega-1\}, \min\{\lambda,
\omega\}]$, also works.

\textbf{(c) Player 2: behavioral $\Rightarrow$ mixed.} A behavioral strategy is $\beta_2 = (\alpha,
\tau, \xi)$: $\Pr(a)$ at the root, $\Pr(t)$ at $h_a$, $\Pr(x)$ at $h_b$. Define the product mixed
strategy
\[
  \sigma_2(r, m_a, m_b) = \Pr\nolimits_\alpha(r)\cdot \Pr\nolimits_\tau(m_a) \cdot \Pr\nolimits_\xi(m_b),
\]
e.g.\ $\sigma_2(a,t,x) = \alpha\tau\xi$. Fix any pure $s_1 = (\ell, m)$ with $\ell \in \{L,R\}$,
$m \in \{W,E\}$. Under $\beta_2$, terminal node $(a, \ell, t)$ has probability $\alpha\tau$. Under
$\sigma_2$ it occurs iff the drawn plan has $r = a$ and $m_a = t$ (any $m_b$):
\[
  \sum_{m_b \in\{x,y\}} \sigma_2(a, t, m_b) = \alpha\tau\xi + \alpha\tau(1-\xi) = \alpha\tau .
\]
The same computation gives $\alpha(1-\tau)$ for $(a, \ell, z)$, $(1-\alpha)\xi$ for $(b, m, x)$ and
$(1-\alpha)(1-\xi)$ for $(b, m, y)$; nodes off $s_1$'s choices have probability 0 under both.
So the two strategies are outcome-equivalent against every $s_1$, hence against every $\sigma_1$.
What makes this work is \emph{perfect recall}: at $h_a$ player~2 remembers she chose $a$, so her
bottom move can be drawn independently of, and after, her root move, and the product of independent
draws reproduces it.

\textbf{(d) Imperfect recall.} Merge $h_a$ and $h_b$ into one information set $h$. Actions at the
nodes of an information set must be the same, so identify $x \equiv t$ and $y \equiv z$. Player~2
now has $S_2' = \{a, b\} \times \{t, z\}$ (4 pure strategies); a behavioral strategy is $(\alpha,
\tau)$; she has \emph{forgotten} her own first move.

\emph{The claim of (c) survives.} The product mixed strategy $\sigma_2(r, m) = \Pr_\alpha(r)
\Pr_\tau(m)$ still reproduces $(\alpha, \tau)$: each play passes through $h$ at most once, so the
probability of a terminal node $(r, \ell, m)$ is $\Pr_\alpha(r)\Pr_\tau(m)$ under both. No example
can break (c) in this game. (Behavioral $\Rightarrow$ mixed fails only when one play can cross the
same information set twice, as in the absent-minded driver, Tadelis p.~142.)

\emph{What fails is the converse} (mixed $\Rightarrow$ behavioral, the other half of Kuhn's
theorem). Take
\[
  \sigma_2 = \tfrac12\,(a, t) + \tfrac12\,(b, z)
\]
(``go left and then play $t$, or go right and then play $z$'': the plan coordinates the bottom move
with the root move she no longer remembers). Against $s_1 = LW$ it gives $\Pr(a,L,t) = \frac12$,
$\Pr(b,W,z) = \frac12$, $\Pr(a,L,z) = 0$. A behavioral $(\alpha, \tau)$ gives
\[
  \Pr(a, L, t) = \alpha\tau, \qquad \Pr(a, L, z) = \alpha(1 - \tau), \qquad
  \Pr(b, W, z) = (1-\alpha)(1-\tau).
\]
$\Pr(a,L,z) = 0$ forces $\alpha = 0$ or $\tau = 1$; the first makes $\Pr(a,L,t) = 0$, the second
makes $\Pr(b,W,z) = 0$, contradicting $\frac12$ in either case. So this mixed strategy has no
outcome-equivalent behavioral strategy. It also matters for payoffs: against $LW$, $\sigma_2$ gives
player~2 $\frac12 \cdot 1 + \frac12 \cdot 0 = \frac12$, and against $RE$ it gives $\frac12\cdot 0 +
\frac12 \cdot 2 = 1$; behavioral strategies cannot replicate this coordination.

\interp Mixed strategies are plans chosen before the game; behavioral strategies are choices made
at each information set. With perfect recall the two are interchangeable, which is why the course
(and Tadelis) can analyze dynamic games with behavioral strategies and define SPE with them
(Lecture~3, slide~34). Without recall a player who randomizes ``in the moment'' cannot correlate
her later move with an earlier one she has forgotten, while a plan made in advance can.

\verify \texttt{check\_B\_L03.py} computes with sympy, for each pure strategy of the opponent and
each of the eight terminal nodes, the probability under the behavioral strategy and under the
product mixed strategy, for player~1 (b) and player~2 (c); all differences simplify to 0. For (d)
it confirms the product construction still works and that the system $\alpha\tau = \frac12$,
$\alpha(1-\tau) = 0$, $(1-\alpha)(1-\tau) = \frac12$ has no solution. Output line:
\texttt{tadelis-7-10: (b), (c) equivalences hold symbolically; (d) mixed->behavioral fails}.
Tadelis prints no answer. \textbf{Note on the statement:} read literally, (d) asks for a failure of
behavioral $\Rightarrow$ mixed, which cannot fail here (each play crosses the merged set once); the
example above shows the failure of the converse, which is the equivalence that perfect recall
protects in this game.

\mistakes
\begin{itemize}
  \item Giving player~2 only $2 \times 2 = 4$ strategies in (a): with perfect recall she has two
  separate bottom information sets, so $2^3 = 8$.
  \item Proving (b) only against one strategy of player~2: equivalence must hold against all of them.
  \item In (d), claiming that the product construction breaks. It does not; the coordination of a
  mixed plan is what behavioral strategies lose.
\end{itemize}

\answer{(a) $S_1 = \{LW, LE, RW, RE\}$; $S_2 = \{a,b\}\times\{t,z\}\times\{x,y\}$ (8 strategies).
(b), (c) The product of the behavioral probabilities is a mixed strategy with the same distribution
over terminal nodes against every opponent strategy. (d) Behavioral $\Rightarrow$ mixed still holds
(each play crosses the merged set once); the equivalence that fails is mixed $\Rightarrow$
behavioral: $\frac12(a,t) + \frac12(b,z)$ has no behavioral equivalent.}
\end{exsol}
